% 2.1 容许表示（底本 p.10--16）
\subsection{$\GL_n$ 的容许表示}

\subsubsection{一般定义}

设 $\mathbf{G}$ 是 $K$ 上的连通约化群，记 $G=\mathbf{G}(K)$。设 $V$ 是 $\mathbb{C}$-向量空间
（下文中 $\mathbb{C}$ 可换成任意代数闭的特征 $0$ 域：事实上我们固定一个同构
$\iota\colon L=\overline{\mathbb{Q}}_l\xrightarrow{\ \sim\ }\mathbb{C}$，此后在笔记中不再区分这些系数域）。

\begin{definition}
设 $(\pi,V)$ 是 $G$ 的表示。
\begin{enumerate}[label=(\arabic*),leftmargin=2.2em]
\item $\pi$ 称为\emph{光滑的}，若每个 $v\in V$ 的稳定化子在 $G$ 中开。
\item $\pi$ 称为\emph{容许的}，若它光滑，且对每个开子群 $H\subset G$，$V^{H}$ 有限维。
\end{enumerate}
\end{definition}

记 $\Repinf_G$ 为光滑表示的范畴，$\Repadm_G$ 为容许表示的范畴。

\begin{remark}
注意 $\pi$ 光滑当且仅当 $V=\bigcup_K V^{K}$，其中 $K$ 遍历 $G$ 的所有紧开子群。
若进一步每个 $V^{K}$ 有限维，则 $\pi$ 容许。
\end{remark}

\begin{remark}
设 $K$ 是 $G$ 的任一紧开子群，$E(K)$ 是光滑不可约表示的同构类集合
（注意光滑性蕴含这些表示必有限维）。限制 $V|_K$ 显然半单：由光滑性，每个 $v\in V$
都落在某个（有限维的）$K$ 的不可约表示里。对每个 $\sigma=(\Sigma,\sigma)\in E(K)$，
记 $V_{\sigma}$ 为 $V$ 的 $\sigma$-等型分量，即 $V$ 中承载该类表示的最小 $K$-不变子空间；
于是 $V=\bigoplus_{\sigma\in E(K)}V_{\sigma}$，且每个 $V_{\sigma}|_K$ 是 $K$ 的表示 $\Sigma$
的直和。$(\pi,V)$ 容许，恰当每个 $V_{\sigma}$ 有限维；此时
\[
V|_K=\bigoplus V_{\sigma}=\bigoplus_{\sigma\in E(K)}\Sigma^{m_{\sigma}},
\]
其中 $m_{\sigma}\in\mathbb{Z}_{\ge 0}$ 是 $\sigma$ 在 $V|_K$ 中的重数。
\end{remark}

\begin{definition}
设 $(\pi,V)\in\Repinf_G$，$V^{*}$ 是 $V$ 的线性对偶（在向量空间范畴中），它按通常方式
带有 $G$ 的作用 $(g\psi)(v)=\psi(g^{-1}v)$。定义
\[
V^{\vee}=\widetilde{V}:=\{\psi\in V^{*}:\ \psi\ \text{的每个稳定化子在}\ G\ \text{中开}\}.
\]
$V^{\vee}=\widetilde{V}$ 称为表示 $V$ 的\emph{反变对偶}（contragradient），也记作
$\pi^{\vee}$ 或 $\widetilde{\pi}$；按定义它是光滑的。
\end{definition}

\begin{remark}
一般地，自然态射 $V\to V^{\vee\vee}$ 不是同构。从某种意义上说，容许性正是为了让这个
问题消失而定义的：
\begin{enumerate}[label=(\arabic*),leftmargin=2.2em]
\item $\pi$ 容许当且仅当 $\pi^{\vee}$ 容许；
\item $\pi$ 容许当且仅当 $V\xrightarrow{\ \sim\ }V^{\vee\vee}$ 是同构；
\item $\pi$ 不可约当且仅当 $\pi^{\vee}$ 不可约；
\item （Gelfand--Kazhdan）当 $G=\GL_n(K)$、$\pi\in\Repinf_G$ 不可约时，$\pi^{\vee}$ 同构于
      表示 $g\mapsto\pi({}^{t}g^{-1})$，这正是「反变对偶」这个名字的来历。
\end{enumerate}
\end{remark}

\subsubsection{Hecke 代数}

\begin{definition}
设 $G$ 如上，把 $G$ 上的 Haar 测度正规化，使 $\mu(K_0)=1$，其中 $K_0\subset G$ 是极大紧子群。
\begin{enumerate}[label=(\arabic*),leftmargin=2.2em]
\item \emph{Hecke 代数} $H(G)$ 是紧支集的局部常值函数 $f\colon G\to\mathbb{C}$ 构成的集合，
      在卷积下（取定 Haar 测度后）成为结合代数：
\[
(f_1*f_2)(h)=\int_G f_1(hg^{-1})f_2(g)\,dg=\int_G f_1(g)f_2(g^{-1}h)\,dg .
\]
\item 设 $K\subset G$ 是紧开子群，定义 $H(G/\!/K)$ 为 $H(G)$ 中在 $K$ 下左右不变的
      函数构成的子代数。
\item 更一般地，设 $Z\subset Z(G)$ 是中心的闭子空间，$\chi\colon Z\to\mathbb{C}$ 是 $Z$ 的
      特征标。有如下双陪集 Hecke 代数的变体：
\[
H_{\chi}(G):=\bigl\{f\colon G\to\mathbb{C}\ \text{局部常值、模}\ Z\ \text{紧支集}:\ 
f(zg)=\chi(z)^{-1}f(g)\ \ \forall z\in Z\bigr\},
\]
其乘法为 $(f_1*f_2)(h)=\int_{G/Z}f_1(x)f_2(x^{-1}h)\,dx$。
\end{enumerate}
\end{definition}

\begin{remark}
\leavevmode
\begin{enumerate}[label=(\arabic*),leftmargin=2.2em]
\item $H(G/\!/K)=\{\varphi\colon K\backslash G/K\to\mathbb{C}\ \text{有限支集}\}$。换言之，
      $H(G,K)$ 由那些满足 $\varphi(kgk')=\varphi(g)$（对一切 $k,k'\in K$）且支集为有限个
      双陪集 $KgK$ 之并的 $\varphi\in H(G)$ 构成。
\item $H(G/\!/K)$ 有双边单位元，即 $e_K:=\frac{1}{\vol(K)}\mathbf{1}_K$，$K$ 上正规化的
      特征函数。
\item 然而 $H(G)$ 不是含幺代数。$H(G)=\bigcup_K H(G/\!/K)$ 是各子代数之并，
      但对不同的 $K$ 单位元 $e_K$ 也不同。
\end{enumerate}
\end{remark}

\subsubsection{$G$ 在 $H(G)$ 上的作用}

$G$ 在 $H(G)$ 上有自然的左作用（相应地，右作用）\en{1}，现描述如下。
\begin{itemize}[leftmargin=1.6em]
\item $G$ 在 $H(G)$ 上的左作用记为 $\lambda_g$，由
\[
\lambda_gf(z)=f(g^{-1}z)\qquad(\forall z\in G)
\]
给出。显然对任意 $f\in H(G)$ 有 $\lambda_gf\in H(G)$，且这确实定义了作用。
\item 同样，右作用 $\rho_g$ 由 $\rho_gf(z)=f(zg)$（$\forall z\in G$）给出。
\end{itemize}

\begin{remark}
我们把 $H(G)$ 定义成 $C^{\infty}_c(G,\mathbb{C})$ 的部分原因，正是要使上面两个 $G$-作用
都成为 $G$ 的光滑表示。
\end{remark}

\subsubsection{Hecke 模}

设 $(\pi,V)$ 是光滑表示，$H(G)$ 通过
\[
f\star v:=\pi(f)(v)=\int_G f(g)\pi(g)v\,dg
\]
作用在 $V$ 上。由上面的注记，这仍是有限和。由于 $V=\bigcup_K V^{K}$，Hecke 作用使 $V$
成为非退化的 $H(G)$-模，即 $H(G)\cdot V=V$（例如 $\frac{1}{\mu(K)}\mathbf{1}_K\star v
=\pi(e_K)v=v$）。

\begin{lemma}
存在自然的函子
\[
\Repinf_G\longrightarrow\{\text{非退化}\ H(G)\text{-模}\},
\]
它诱导出范畴等价。
\end{lemma}

在继续之前，我们固定一些记号并记录几条简单观察。

\begin{enumerate}[label=(\arabic*),leftmargin=2.2em]
\item 对任意紧开子群 $K\subset G$，令 $e_{gK}:=\frac{1}{\mu(K)}\mathbf{1}_{gK}$ 为 $gK$ 上
      正规化的特征函数。注意当 $g=\mathrm{id}$ 时，它与 $e_K$ 的记号一致。
\item 直接验证可知：对 $f\in H(G)$，$e_K*f=f$ 当且仅当对一切 $k\in K$ 有 $f(kg)=f(g)$。
      特别地，对任意紧开子群 $K'\subset K$，有 $e_K=e_{K'}*e_K$。
\item 仍设 $K'\subset K$ 如上，$g\in G$ 任意。我们断言 $e_{gK}=e_{gK'}*e_K$。
      由定义，
\[
e_{gK'}*e_K(z)=\int_G e_{gK'}(x)\,e_K(x^{-1}z)\,dx,
\]
它非零当且仅当 $z\in xK$、其中 $x\in gK'$，即当且仅当 $z\in gK$。对这样的 $z$，
$e_{gK'}*e_K(z)$ 的值恰为
$\frac{1}{\mu(K)}\mathbf{1}_{K'}\cdot\mu(gK'\cap gK)=\frac{1}{\mu(K)}$。
此处关键在于第二个特征函数支撑在一个\emph{子群} $K$ 上（而非它的某个平移 $gK$ 上）。
\item 在上面描述的函子之下，$(\pi,V)$ 的子表示对应于 $H(G)$ 作用下的子模。
\end{enumerate}

\begin{proof}[引理 2.1 的证明]
我们已经构造了从 $\Repinf_G$ 到非退化 Hecke 模的函子；展开定义可以验证它全忠实。
现构造反方向的函子。设 $M$ 是非退化 Hecke 模。首先断言：对任意 $m\in M$，存在紧开子群
$K\subset G$ 使 $e_K\star m=m$。事实上，由非退化性可写 $m=\sum f_i\star m_i$（有限和，
$f_i\in H(G)$，$m_i\in M$）。取 $K$ 足够小，使 $K\subset G$ 在 $H(G)$ 上的左右作用都稳定
每个 $f_i$（见 2.1.3 小节），即 $f_i\in H(G/\!/K)$。于是
\[
e_K\star m=e_K\star\Bigl(\sum f_i\star m_i\Bigr)
=\sum\bigl((e_K*f_i)\star m_i\bigr)=m .
\]
取这样的紧开子群 $K$，定义
\[
g\cdot m=\pi(g)m:=e_{gK}\star m .
\]
首先这定义良好，即不依赖于 $K$ 的选取（只要它足够小使得 $e_K\star m=m$）。这很直接：
设 $K'\subset K$，则
\[
e_{gK'}\star m=e_{gK'}\star(e_K\star m)=(e_{gK'}*e_K)\star m=e_{gK}\star m,
\]
最后一步用的是上面的观察 (3)。

接着要证明这确实定义了作用，这一点其实颇为微妙（但光滑性是直接的）。
设 $g_1,g_2\in G$、$m\in M$，我们要选取 $K$ 与 $K'$ 使
\[
e_{g_1g_2K}\star m=e_{g_1K'}\star e_{g_2K}\star m
=(e_{g_1K'}*e_{g_2K})\star m .
\]
$K$ 的选取如前，使得 $e_K\star m=m$。现在令 $m':=g_2\cdot m$，断言可取 $K'$ 足够小，
使得 $e_{g_2K}\in H(G/\!/K')$ 且 $e_{K'}\star m'=m'$。为此注意 $e_{g_2K}\in H(G/\!/K')$
意即 $K'g_2K=g_2K$；先取 $K'=g_2Kg_2^{-1}\cap K$\en{2}，必要时再缩小以稳定 $m'$。
由此可直接验证
\[
e_{g_1K'}*e_{g_2K}=e_{g_1g_2K},
\]
从而确有 $e_{g_1g_2K}\star m=(e_{g_1K'}*e_{g_2K})\star m$；再把 $K$ 与 $K'$ 换成它们的交，
即得 $(g_1g_2)\cdot m=g_1\cdot(g_2\cdot m)$。
\end{proof}

\begin{remark}
同一构造还有另一种（更简洁的）做法：考虑 $H(G)$-模的自然映射
\[
H(G)\otimes_{H(G)}M\longrightarrow M,
\]
其中把 $H(G)$ 看作带右作用在其自身上的模。该映射其实是同构（尽管 $H(G)$ 不含单位元），
这来自非退化假设。于是通过同构 $H(G)\otimes M\cong M$、并让 $G$ 左作用在 $H(G)$ 上，
就得到 $M$ 上的 $G$-作用。注意这恰好对应
$g\cdot m=g\cdot(e_K\otimes m)=e_{gK}\otimes m$。
\end{remark}

\subsubsection{}

现在固定 $K$，则上面的函子限制为从 $\Repadm_G$ 到有限维 $H(G/\!/K)$-模范畴的函子。
在本书中，「不可约模」与「单模」作为同义词使用。

\begin{lemma}[Casselman]
固定紧开子群 $K\subset G$。
\begin{enumerate}[label=(\arabic*),leftmargin=2.2em]
\item 函子
\[
\bigl\{\pi\in\Rep^{f.l.}_{G}{}^{\mathrm{adm}}:\ 
\forall W\ \text{不可约子商},\ W^{K}\neq 0\bigr\}\xrightarrow{\ \sim\ }H(G/\!/K)\text{-}\mathrm{mod}^{f.d.},
\qquad V\longmapsto V^{K}
\]
是从「所有不可约子商都含 $K$-不动向量的有限长容许表示」范畴到有限维
$H(G/\!/K)$-模范畴的等价。
\item 若 $V\in\Repinf_G$ 不可约且 $V^{K}\neq 0$，则 $V^{K}$ 是不可约 $H(G/\!/K)$-模。
      此外，下述两类对象之间存在双射：
\begin{enumerate}[label=(\alph*),leftmargin=2.0em]
\item 满足 $V^{K}\neq 0$ 的不可约光滑表示 $(\pi,V)$ 的等价类；
\item 不可约 $H(G/\!/K)$-模的同构类。
\end{enumerate}
\end{enumerate}
\end{lemma}

\begin{remark}
容易证明：若 $V\in\Repinf_G$ 不可约，则 $V^{K}$ 或是 $0$、或是不可约 $H(G)$-模。
但反过来则是一个非平凡的陈述。
\end{remark}

作为推论，我们有

\begin{corollary}
若 $V\in\Repinf_G$ 不可约且 $H(G/\!/K)$ 交换，则或者 $V^{K}=0$，或者 $V^{K}$ 是一维的。
\end{corollary}

\subsubsection{球面 Hecke 代数}

在我们的设定里 $G=\GL_n(K)$。若 $C=\GL_n(\mathcal{O}_K)$ 是 $G$ 的极大紧开子群，
则 $H(G/\!/C)$ 交换；事实上
\[
H(G/\!/C)\xrightarrow{\ \sim\ }\mathbb{C}[t_1^{\pm 1},\dots,t_n^{\pm 1}]^{S_n}.
\]
更一般地，对 $G=\mathbf{G}(K)$（$\mathbf{G}$ 分裂约化，固定 $\mathbf{T}\subset\mathbf{B}\subset\mathbf{G}$）
与超特殊子群 $C\subset G$，经典的 Satake 同构说
\[
H(G/\!/C)\xrightarrow{\ \sim\ }(\mathcal{O}_{\widehat{T}})^{W},
\]
其中 $\widehat{T}=\Spec\mathbb{C}[\Lambda]$，$\Lambda$ 是余权格，$W$ 是 Weyl 群。
特别地，球面 Hecke 代数交换。

\begin{remark}
在 $\GL_n$ 情形，$H$ 的交换性也可以从 Gelfand 技巧看出。注意
\[
C\backslash G/C\cong G\backslash\bigl(G/C\times G/C\bigr),
\qquad [g]\longmapsto(1,[g]).
\]
而 $G/C$ 参数化 $K^{n}$ 中的格。也就是说可以等同
\[
C\backslash G/C\ \longleftrightarrow\ G\backslash(\mathcal{L}^{n}\times\mathcal{L}^{n})
=\{\text{格的相对位置}\},
\]
其中用 $\mathcal{L}^{n}$ 记 $K^{n}$ 中的格。从这个观点看，
\[
(\varphi*\psi)(\Lambda,\Lambda')=\sum_{\Lambda''}\varphi(\Lambda,\Lambda'')\psi(\Lambda'',\Lambda').
\]
注意此和有限，因为支集紧。Gelfand 的想法是：$(\Lambda,\Lambda')$ 的相对位置与
$((\Lambda')^{\vee},\Lambda^{\vee})$ 相同\en{3}，因此
\begin{multline*}
(\varphi*\psi)(\Lambda,\Lambda')
=\sum_{\Lambda''}\varphi(\Lambda,\Lambda'')\psi(\Lambda'',\Lambda')
=\sum_{\Lambda''}\varphi\bigl((\Lambda'')^{\vee},\Lambda^{\vee}\bigr)
\psi\bigl((\Lambda')^{\vee},(\Lambda'')^{\vee}\bigr)\\
=\sum_{\Lambda''}\varphi(\Lambda'',\Lambda^{\vee})\psi\bigl((\Lambda')^{\vee},\Lambda''\bigr)
=(\psi*\varphi)\bigl((\Lambda')^{\vee},\Lambda^{\vee}\bigr)=(\psi*\varphi)(\Lambda,\Lambda').
\end{multline*}
这个证明几乎是同义反复。
\end{remark}

\begin{corollarydefinition}
设 $(\pi,V)$ 是 $G$ 的不可约容许表示，$C=\GL_n(\mathcal{O}_K)$ 是 $G$ 的极大紧子群，
则 $V^{C}$ 或为 $0$、或为一维。若 $V^{C}\neq 0$（底本此处误作 $V^{\pi}\neq 0$，见译注），
则称 $V$ 是\emph{非分歧的}，即 $V^{C}$ 一维。
\end{corollarydefinition}

\begin{yizhu}[底本的笔误：$V^{\pi}$ 应为 $V^{C}$]
底本 2.1.6 末写作「$V$ is called unramified if $V^{\pi}\neq 0$, i.e., $V^{\pi}$ is
1-dimensional」。此处 $\pi$ 应是 $C$：全文用 $V^{C}$ 记 $C$-不动向量，
而 $V^{\pi}$ 没有意义（$\pi$ 是表示本身）。译文按 $V^{C}$ 处理，并在正文括注原文。
\end{yizhu}

\subsubsection{Schur 引理}

\begin{lemma}
设 $(\pi,V)\in\Repinf_G$ 是 $G$ 的光滑不可约表示，则 $V$ 的每个 $G$-自同态都是数乘。
\end{lemma}

与 $V$ 有限维的情形不同，Schur 引理在这里是非平凡的陈述。当
$\pi\in\Repadm_G$ 不可约且容许\en{4}时，可以借引理 2.2 回到有限维模的情形：
取 $C\subset G$ 足够小使 $V^{C}\neq 0$，则 $V^{C}$ 的 $H(G/\!/C)$-线性自同态只有数乘。

\begin{proof}[略证]
在我们的设定 $G=\mathbf{G}(K)$ 下，注意 $G$ 有可数基，且对每个紧开子群 $K\subset G$，
$H(G/\!/K)$ 在 $\mathbb{C}$ 上可数维，故 $H(G)$ 在 $\mathbb{C}$ 上可数维。
取 $V$ 中非零向量 $v$，由 $(\pi,V)$ 不可约，映射 $H(G)\to V$、$f\mapsto f\star v$ 是满射，
故 $V$ 可数维。再考虑环 $\End_G(V)$：仍由 $(\pi,V)$ 的不可约性，
映射 $\End_G(V)\to V$、$\mu\mapsto\mu(v)$ 是单射。于是 $\End_G(V)$ 是 $\mathbb{C}$ 上可数维的
除环；但不可数代数闭域 $F$ 上可数维的除环必为 $F$ 本身（证明：若 $D\neq F$，
取 $x\in D\setminus F$，则 $x$ 在 $F$ 上超越——因为 $F$ 代数闭；而元素族
$\{\frac{1}{x-\alpha}:\alpha\in F\}$ 不可数，且可验证其在 $F$ 上线性无关。证毕），
于是引理成立。
\end{proof}

\subsubsection{光滑不可约表示}

\begin{theorem}
设 $(\pi,V)$ 是 $G$ 的光滑不可约表示，则 $\pi$ 容许。
\end{theorem}

证明用 Jacquet 函子把 $\pi$ 嵌入到容许（超尖点）表示的抛物诱导中。
我们略去证明，请读者参阅 Jacquet 的处理或 [Bernstein--Zelevinsky]。

\begin{lemma}
$G=\GL_n(K)$ 的每个光滑不可约表示或是一维的、或是无限维的。
\end{lemma}

\begin{proof}
这来自 $G$ 的结构。设 $\pi\colon G\to\GL_n(\mathbb{C})$ 光滑，则
$\ker(\pi)\triangleleft G$ 是开子群，因而包含某个 $1+\varpi^{m}\mathfrak{M}_n(\mathcal{O}_K)
\cong 1+\varpi^{m}M_n(\mathcal{O}_K)$，从而包含 $1+\varpi^{m}\mathfrak{n}$，
即对角元为 $1$、其余元素落在 $\varpi^{m}\mathcal{O}_K$ 中的上三角矩阵全体。
由共轭可知幂幺上三角矩阵 $U\subset\ker\pi$。最后由正规性知 $\SL_n(K)\subset\ker\pi$。
这样我们就证明了引理，并得到更多：任何这样的有限维（从而一维）不可约表示 $\pi$
都通过行列式映射分解，即形如 $\chi\circ\det$，其中 $\chi$ 是 $K^{\times}$ 的某个拟特征标。
\end{proof}

\subsubsection{Harish-Chandra 特征标}\en{5}

回忆：给定光滑表示 $(\pi,V)$，Hecke 代数 $H(G)$ 通过
\[
\pi(f)(v)=\int_G f(g)\pi(g)v\,dg,\qquad f\in H(G)=C^{\infty}_c(G)
\]
作用在 $V$ 上。在「$G$ 的光滑表示 $\leftrightarrow$ 非退化 $H(G)$-模」的等价之下，
容许表示对应于那些「对一切 $f$，$\pi(f)$ 是有限秩算子」的 $H(G)$-模。
这使我们可以定义 $\pi(f)$ 的迹，从而定义分布特征标
\[
\Theta_{\pi}\colon H(G)\to\mathbb{C},\qquad f\longmapsto\Theta_{\pi}(f):=\Tr\pi(f).
\]

\begin{lemma}
若 $\{\pi_1,\dots,\pi_n\}$ 是 $G$ 的两两不同构的不可约容许表示，则
$\{\Theta_{\pi_1},\dots,\Theta_{\pi_n}\}$ 线性无关。
\end{lemma}

\begin{remark}
计算分布特征标 $\Theta_{\pi}$（$(\pi,V)$ 给定）有一个有用的办法如下。
取 $G$ 的一列可数基 $K_m$（例如 $K_i=1+\varpi^{i}K_0$，$K_0=\GL_2(\mathcal{O}_K)$\en{原注}），
令 $V_i:=V^{K_i}$，则 $V_0\subset V_1\subset\cdots\subset V=\bigcup_iV_i$。对每个 $i$ 与
每个 $g\in G$，考虑复合
\[
\pi(e_{K_i})\circ\pi(g)\circ\pi(e_{K_i})\colon V_i\longrightarrow V_i,
\]
记其迹为 $\xi_{\pi,i}(g)$。对 $f\in H(G)$，存在某个 $i$（依赖于 $f$）使 $f\in H(G/\!/K_i)$；
于是对每个 $m\ge i$ 有
\[
\Theta_{\pi}(f)=\int_G\xi_{\pi,m}(g)f(g)\,dg,
\]
且右端在 $m$ 足够大时不依赖于 $m$。
\end{remark}

我们陈述 Harish-Chandra 的一个深刻定理（事实上，当 $K$ 特征为 $0$ 时这是 Harish-Chandra
的定理，正特征时是 Lemaire 的定理）：$\Theta_{\pi}$ 可由 Harish-Chandra 特征标
$\chi_{\pi}$ 表示，后者在正则半单元 $G^{\mathrm{r.s.}}$ 上局部常值。先建立记号：
\begin{enumerate}[label=(\arabic*),leftmargin=2.2em]
\item 记 $G^{\mathrm{r.s.}}$ 为 $\GL_n(K)$ 中正则半单元的集合。（回忆 $g\in\GL_n(K)$
      正则半单，指其特征多项式 $P_g(X)\in k[X]$ 在 $K$ 上无重根。）
\item 记 $G^{\mathrm{r.e.}}\subset G^{\mathrm{r.s.}}$ 为正则椭圆元的子集。（回忆
      $g\in G^{\mathrm{r.s.}}$ 正则椭圆，指此外其特征多项式 $P_g$ 不可约。）
\end{enumerate}

\begin{theorem}[Harish-Chandra]
存在 $G$ 上的局部可积函数 $\chi_{\pi}$，它在正则半单集 $G^{\mathrm{r.s.}}$ 上局部常值、
在 $G^{\mathrm{r.s.}}$ 的补集上为 $0$，使得
\[
\Theta_{\pi}(f)=\int_G f(g)\chi_{\pi}(g)\,dg
\]
对一切 $f\in H(G)$ 成立。
\end{theorem}

\begin{remark}
$\chi_{\pi}$ 可以看作特征标 $\chi_{\pi}\colon\widetilde{G}\to\mathbb{C}$，其中
$\widetilde{G}=G^{\mathrm{r.s.}}/\sim$ 是 $G$ 中正则半单元的共轭类集合。
\end{remark}

\begin{remark}
我们记录一个在讨论 Jacquet--Langlands 对应时有用的定义。
\begin{enumerate}[label=(\arabic*),leftmargin=2.2em]
\item 光滑容许表示 $(\pi,V)$ 称为\emph{椭圆的}，若其 Harish-Chandra 特征标
      $\chi_{\pi}$ 在 $G^{\mathrm{r.e.}}$ 上不恒为 $0$。
\item 每个本质平方可积表示都是椭圆的。
\end{enumerate}
\end{remark}
